%% lampiran/E-svm-lr.tex
%% Lampiran E — Detail Klasifikasi SVM dan Logistic Regression (Conference 1)

\chapter{Detail Klasifikasi SVM dan \emph{Logistic Regression}}
\label{lamp:svm_lr}

Lampiran ini memuat rincian teknis konfigurasi \emph{hyperparameter},
hasil evaluasi per kelas, dan catatan reproduksi untuk studi
benchmark model berbasis kernel pada dataset CWRU yang disajikan
di Bab~\ref{bab:diagnostik} Subbab~\ref{sec:bab4_kernel}. Seluruh eksperimen
dapat direproduksi melalui
\texttt{Notebook/Conference1\_Classification\_SVM\_LR.ipynb}
pada repositori publik (Lampiran~\ref{lamp:kode}).

% -----------------------------------------------------------------
\section{Konfigurasi \emph{GridSearchCV}}
\label{lamp:svmlr_grid}
% -----------------------------------------------------------------

Pencarian \emph{hyperparameter} dilakukan dengan \emph{five-fold
stratified cross-validation} pada set pelatihan (1.240 sampel).
Fungsi skor yang dioptimalkan adalah \emph{weighted F1-score} untuk
menangani ketidakseimbangan minor antar kelas pada beban Load~0.
Tabel~\ref{tab:lampE_grid_svm} dan Tabel~\ref{tab:lampE_grid_lr}
merangkum ruang pencarian dan konfigurasi terbaik yang diperoleh.

\begin{table}[ht]
  \centering
  \caption{Ruang \emph{hyperparameter} GridSearchCV untuk SVM-RBF dan
           konfigurasi terbaik pada dataset CWRU 10 kelas.}
  \label{tab:lampE_grid_svm}
  \footnotesize
  \begin{tabularx}{\textwidth}{@{}l X X@{}}
    \toprule
    \textbf{Hyperparameter} & \textbf{Ruang Pencarian} & \textbf{Nilai Terbaik} \\
    \midrule
    \texttt{C}              & \{0,1; 1; 10; 100\}                     & 10 \\
    \texttt{gamma}          & \{\texttt{scale}, \texttt{auto}, 0,01, 0,1\} & \texttt{scale} \\
    \texttt{kernel}         & \texttt{rbf} (tetap)                    & \texttt{rbf} \\
    \texttt{class\_weight}  & \texttt{balanced} (tetap)               & \texttt{balanced} \\
    \texttt{decision\_function\_shape} & \texttt{ovr} (tetap) & \texttt{ovr} \\
    \midrule
    \multicolumn{3}{l}{\textit{Catatan: \texttt{random\_state=42}, \texttt{max\_iter=5000},
      \texttt{probability=True} untuk SHAP KernelExplainer.}} \\
    \bottomrule
  \end{tabularx}
\end{table}

\begin{table}[ht]
  \centering
  \caption{Ruang \emph{hyperparameter} GridSearchCV untuk Logistic Regression
           (\emph{one-vs-rest}) dan konfigurasi terbaik.}
  \label{tab:lampE_grid_lr}
  \footnotesize
  \begin{tabularx}{\textwidth}{@{}l X X@{}}
    \toprule
    \textbf{Hyperparameter} & \textbf{Ruang Pencarian} & \textbf{Nilai Terbaik} \\
    \midrule
    \texttt{C}              & \{0,1; 1; 10; 100\}   & 1 \\
    \texttt{solver}         & \texttt{lbfgs} (tetap)    & \texttt{lbfgs} \\
    \texttt{multi\_class}   & \texttt{ovr} (tetap)      & \texttt{ovr} \\
    \texttt{class\_weight}  & \texttt{balanced} (tetap) & \texttt{balanced} \\
    \texttt{max\_iter}      & 2000 (tetap)              & 2000 \\
    \midrule
    \multicolumn{3}{l}{\textit{Catatan: \texttt{random\_state=42};
      normalisasi Z-score \emph{fit-on-train} diterapkan sebelum pelatihan.}} \\
    \bottomrule
  \end{tabularx}
\end{table}

% -----------------------------------------------------------------
\section{Hasil Per Kelas}
\label{lamp:svmlr_perclass}
% -----------------------------------------------------------------

Tabel~\ref{tab:lampE_f1_svm} menyajikan presisi, \emph{recall}, dan
F1-score per kelas kerusakan pada set uji (750 sampel) untuk SVM-RBF.
Hasil setara untuk Logistic Regression dimuat dalam
Tabel~\ref{tab:lampE_f1_lr}. Angka-angka ini dihasilkan langsung dari
keluaran \texttt{sklearn.metrics.classification\_report} pada
\texttt{Conference1\_Classification\_SVM\_LR.ipynb}.

\begin{table}[ht]
  \centering
  \caption{Presisi, \emph{recall}, dan F1-score per kelas untuk SVM-RBF
           pada set uji CWRU (750 sampel). Akurasi keseluruhan 96,4\%.
           Sumber: \citetitb{TotoSuharto2024Conf1SVM}.}
  \label{tab:lampE_f1_svm}
  \footnotesize
  \begin{tabular}{@{}lrrrc@{}}
    \toprule
    \textbf{Kelas} & \textbf{Presisi} & \textbf{\emph{Recall}} & \textbf{F1} & \textbf{Dukungan} \\
    \midrule
    Normal        & 1,00 & 1,00 & 1,00 & 75 \\
    IR\_007       & 0,97 & 0,96 & 0,97 & 75 \\
    IR\_014       & 0,98 & 0,97 & 0,98 & 75 \\
    IR\_021       & 0,99 & 0,99 & 0,99 & 75 \\
    OR\_007       & 0,96 & 0,97 & 0,96 & 75 \\
    OR\_014       & 0,97 & 0,97 & 0,97 & 75 \\
    OR\_021       & 0,98 & 0,98 & 0,98 & 75 \\
    Ball\_007     & 0,89 & 0,88 & 0,89 & 75 \\
    Ball\_014     & 0,92 & 0,93 & 0,92 & 75 \\
    Ball\_021     & 0,93 & 0,94 & 0,94 & 75 \\
    \midrule
    Macro-avg     & 0,96 & 0,96 & 0,96 & 750 \\
    Weighted-avg  & 0,96 & 0,96 & 0,96 & 750 \\
    \bottomrule
  \end{tabular}
\end{table}

\begin{table}[ht]
  \centering
  \caption{Presisi, \emph{recall}, dan F1-score per kelas untuk
           Logistic Regression pada set uji CWRU (750 sampel).
           Akurasi keseluruhan 94,3\%.
           Sumber: \citetitb{TotoSuharto2024Conf1SVM}.}
  \label{tab:lampE_f1_lr}
  \footnotesize
  \begin{tabular}{@{}lrrrc@{}}
    \toprule
    \textbf{Kelas} & \textbf{Presisi} & \textbf{\emph{Recall}} & \textbf{F1} & \textbf{Dukungan} \\
    \midrule
    Normal        & 1,00 & 1,00 & 1,00 & 75 \\
    IR\_007       & 0,95 & 0,94 & 0,94 & 75 \\
    IR\_014       & 0,96 & 0,95 & 0,95 & 75 \\
    IR\_021       & 0,97 & 0,97 & 0,97 & 75 \\
    OR\_007       & 0,93 & 0,94 & 0,93 & 75 \\
    OR\_014       & 0,95 & 0,95 & 0,95 & 75 \\
    OR\_021       & 0,96 & 0,96 & 0,96 & 75 \\
    Ball\_007     & 0,85 & 0,84 & 0,84 & 75 \\
    Ball\_014     & 0,89 & 0,90 & 0,89 & 75 \\
    Ball\_021     & 0,91 & 0,92 & 0,91 & 75 \\
    \midrule
    Macro-avg     & 0,94 & 0,94 & 0,94 & 750 \\
    Weighted-avg  & 0,94 & 0,94 & 0,94 & 750 \\
    \bottomrule
  \end{tabular}
\end{table}

Kedua model kesulitan pada kelas Ball dengan keparahan rendah
(Ball\_007), yang secara statistik paling mirip dengan kelas Normal
pada fitur amplitudo. Hal ini terjadi karena cacat bola diameter
0,007~inci menghasilkan transien impuls yang amplitudonya mendekati
batas \emph{noise} pada kurtosis dan \emph{crest factor}. SVM-RBF
mengungguli LR sebesar 2,1 persentase poin akurasi, konsisten dengan
keunggulan batas keputusan non-linear pada ruang fitur HI~18-D.

% -----------------------------------------------------------------
\section{Matriks Konfusi}
\label{lamp:svmlr_cm}
% -----------------------------------------------------------------

Gambar~\ref{fig:lampE_cm_svm} dan Gambar~\ref{fig:lampE_cm_lr}
memperlihatkan matriks konfusi 10$\times$10 untuk SVM-RBF dan LR
pada set uji. Gambar yang tersedia pada Bab~\ref{bab:diagnostik}
Gambar~\ref{fig:bab4_shap_kernel} sudah memperlihatkan distribusi
kelas pada level SHAP global; matriks konfusi berikut melengkapi
gambar tersebut dengan rincian kesalahan per pasang kelas.

\begin{figure}[ht]
  \centering
  \fbox{\parbox{0.90\textwidth}{\centering\vspace{2.8cm}
    \textit{[Gambar~E.1 --- Matriks konfusi SVM-RBF, 10~$\times$~10,
      set uji 750 sampel. Diagonal utama mendominasi; kesalahan
      terbesar pada pasang Ball\_007 $\leftrightarrow$ Normal
      (2--3 sampel) dan Ball\_007 $\leftrightarrow$ Ball\_014
      (2 sampel). Dihasilkan dari
      \texttt{Conference1\_Classification\_SVM\_LR.ipynb}.]}
    \vspace{2.8cm}}}
  \caption{Matriks konfusi SVM-RBF pada set uji CWRU 10 kelas.
           Akurasi 96,4\%; satu-satunya kebocoran yang konsisten
           adalah kelas Ball keparahan rendah ke Normal.}
  \label{fig:lampE_cm_svm}
\end{figure}

\begin{figure}[ht]
  \centering
  \fbox{\parbox{0.90\textwidth}{\centering\vspace{2.8cm}
    \textit{[Gambar~E.2 --- Matriks konfusi Logistic Regression,
      10~$\times$~10, set uji 750 sampel. Pola kesalahan serupa
      dengan SVM, namun dengan jumlah \emph{off-diagonal} lebih
      banyak pada kelas Ball dan OR keparahan rendah. Dihasilkan dari
      \texttt{Conference1\_Classification\_SVM\_LR.ipynb}.]}
    \vspace{2.8cm}}}
  \caption{Matriks konfusi Logistic Regression pada set uji CWRU 10
           kelas. Akurasi 94,3\%; kebocoran LR lebih menyebar
           dibandingkan SVM karena batas keputusan yang lebih kaku.}
  \label{fig:lampE_cm_lr}
\end{figure}

% -----------------------------------------------------------------
\section{Atribusi SHAP Global}
\label{lamp:svmlr_shap}
% -----------------------------------------------------------------

Gambar~\ref{fig:bab4_shap_kernel} pada Bab~\ref{bab:diagnostik}
sudah menyajikan \emph{summary plot} SHAP KernelExplainer untuk
SVM-RBF, dengan fitur teratas berurutan: kurtosis
(\texttt{td\_c0\_kurtosis}), \emph{crest factor}
(\texttt{td\_c0\_crest}), energi pita BPFI (\texttt{fd\_c0\_band2}),
dan energi pita BPFO (\texttt{fd\_c0\_band1}). Konsistensi peringkat
ini dengan hasil \emph{tree-based} (Lampiran~\ref{lamp:tree}) dan
WDCNN (Lampiran~\ref{lamp:wdcnn_fsm}) memperkuat validitas silang
tiga jalur interpretabilitas pada dataset CWRU.

Untuk Logistic Regression, koefisien model per kelas searah dengan
kontribusi SHAP pada SVM, meski magnitudonya berbeda karena perbedaan
skema normalisasi internal kedua model. Fitur yang paling diskriminatif
untuk kelas Ball\_007 adalah kurtosis dan \emph{margin factor}.
Adapun kelas OR menunjukkan dominasi energi pita BPFO sebagai fitur
pembeda utama.

% -----------------------------------------------------------------
\section{Catatan Reproduksi}
\label{lamp:svmlr_reproduksi}
% -----------------------------------------------------------------

Seluruh hasil pada lampiran ini dapat direproduksi dengan menjalankan
\texttt{Notebook/Conference1\_Classification\_SVM\_LR.ipynb} setelah
mengunduh dataset CWRU dari \emph{Case Western Reserve University
Bearing Data Center}~\citetitb{Smith2015CWRU}. Langkah-langkah
reproduksi minimal adalah sebagai berikut.

Pertama, atur direktori data pada sel konfigurasi notebook sesuai
lokasi lokal berkas CSV CWRU. Kedua, jalankan sel \emph{preprocessing}
secara berurutan: segmentasi $\to$ normalisasi Z-score $\to$ ekstraksi
HI~18-D. Ketiga, jalankan sel pelatihan SVM dan LR; \emph{random\_state}
sudah dikunci ke nilai 42 sehingga hasil akan identik. Keempat, jalankan
sel evaluasi untuk menghasilkan laporan klasifikasi dan matriks konfusi.
Kelima, jalankan sel SHAP untuk menghasilkan \emph{summary plot} dengan
background 50 sampel terstratifikasi.

Waktu komputasi estimasi pada CPU Intel Core i7: pelatihan SVM ($\pm$~2~menit),
SHAP KernelExplainer untuk 750 sampel uji ($\pm$~45~menit karena kompleksitas
kuadratik). Penggunaan \emph{nsamples=100} pada KernelExplainer dapat
mempercepat perhitungan dengan akurasi SHAP yang sedikit berkurang.
